(1) 取平面波 \(\psi=e^{i(kx-\omega t)}\) 代入 Klein--Gordon 方程，得到 \(\omega^2=c^2k^2+m^2c^4/\hbar^2\)。正能支为 \(E=\hbar\omega=\sqrt{(\hbar ck)^2+m^2c^4}\)，动量 \(p=\hbar k\)。当 \(k=0\) 时 \(E=mc^2\)，是粒子的静能，而不是非相对论零能基准。

(2) 相速度 \(v_p=\omega/k=\sqrt{c^2+m^2c^4/(\hbar^2k^2)}>c\)，群速度 \(v_g=d\omega/dk=c^2k/\omega=pc^2/E<c\)。两者满足 \(v_pv_g=c^2\)，所以超光速的相位峰并不传递信息；足够窄的正能波包以群速度输运其中心。不能把单频无穷平面波相位传播当作可调制信号前沿。

(3) 二阶色散系数为 \(\omega''(k)=d(c^2k/\omega)/dk=c^2(\omega^2-c^2k^2)/\omega^3=m^2c^6/(\hbar^2\omega^3)>0\)。对中心波数附近作 Taylor 展开，\(\omega(k+\delta k)=\omega(k)+v_g\delta k+\omega''\delta k^2/2+\cdots\)。不同波数群速不同，窄波包随时间发生展宽；在高能质量项相对小的极限，\(\omega''\to0\) 而近似无色散。

(4) 设 \(\epsilon=\hbar k/(mc)\ll1\)，则 \(E=mc^2\sqrt{1+\epsilon^2}=mc^2+(\hbar k)^2/(2m)-(\hbar k)^4/(8m^3c^2)+\cdots\)。扣除静能 \(mc^2\) 后首项就是 Schrödinger 非相对论自由粒子动能 \(p^2/(2m)\)。群速度展开 \(v_g=p/m+O(p^3/(m^3c^2))\)，恢复 Newtonian 速度；静能相位虽快，但不改变扣除常数后的慢速包络动力学。
